\documentclass[12pt]{article}
\usepackage[T1]{fontenc}
\usepackage[utf8]{inputenc}
\usepackage{lmodern}
\usepackage{textcomp}
\usepackage{enumitem}
\usepackage{geometry}
\geometry{margin=1in}
\begin{document}
\begin{center}
Answer Key - Psychology - Graduate Level Assessment 8 \
\small (Answers are ordered exactly as in the question paper.)
\end{center}

\section*{Multiple Choice}
\begin{enumerate}[label=\arabic*.]
\item \textbf{Answer:} D
\\ \textit{Options:} A) occipital lobe; B) temporal lobe; C) parietal lobe; D) cerebral cortex; E) limbic system
\\ \textit{Note:} The cerebral cortex is the least developed area of the brain at birth because it undergoes significant growth and maturation during the early years of life, while other brain structures are more fully developed at birth to support essential survival functions.
\item \textbf{Answer:} C
\\ \textit{Options:} A) personal opinions and assumptions about the client; B) the psychologist's preferred treatment plan for the client; C) interpretations of the data and limiting circumstances involving the test administration; D) only the results that support the initial hypothesis; E) only positive results and findings from the test
\\ \textit{Note:} The correct answer emphasizes the ethical responsibility of psychologists to provide a comprehensive understanding of test results by including both interpretations and any limiting circumstances that may affect the validity of the assessments.
\item \textbf{Answer:} A
\\ \textit{Options:} A) Physical deviation has no impact on adolescents; B) Physical deviation leads to superior athletic performance in adolescents; C) Physical deviation causes a heightened sense of physical well-being in adolescents; D) Physical deviation diminishes the importance of peer relationships in adolescents; E) Physical deviation causes excitement in adolescents
\\ \textit{Note:} The assertion that physical deviation has no impact on adolescents is supported by research indicating that while some individuals may experience temporary psychological effects, many adolescents demonstrate resilience and adapt to their physical differences without significant long-term psychological consequences.
\item \textbf{Answer:} A
\\ \textit{Options:} A) a decrease in the number of psychiatric hospitals.; B) an increase in the number of psychiatric patients.; C) an increase in the number of trained mental health professionals.; D) the development of new psychotherapy techniques.; E) a decrease in government funding.
\\ \textit{Note:} The correct answer is supported by the fact that the mid-1900 s saw significant reductions in the number of psychiatric hospitals due to shifts in mental health policy, which emphasized community-based care and the belief that many patients could be treated effectively outside institutional settings.
\item \textbf{Answer:} A
\\ \textit{Options:} A) Decreased blood pressure, heightened sensory perception, faster reaction times; B) Strengthened cognitive abilities, longer lifespan, better pain tolerance; C) Enhanced problem-solving skills, increased emotional stability, improved coordination; D) Sharper focus, accelerated learning ability, increased creativity; E) Improved cardiovascular health, reduced inflammation, increased energy levels
\\ \textit{Note:} The correct answer highlights that sleep deprivation can lead to physiological changes such as decreased blood pressure, as well as cognitive and sensory enhancements, including heightened sensory perception and faster reaction times, which are well-documented effects in psychological research.
\end{enumerate}

\section*{True / False}
\begin{enumerate}[label=\arabic*.]
\setcounter{enumi}{5}
\item \textbf{Answer:} False
\\ \textit{Options:} A) True; B) False
\\ \textit{Note:} The statement is false because the three attributes of color are hue, saturation, and brightness, not intensity, shade, and tone, which do not accurately describe the fundamental characteristics that differentiate colors.
\item \textbf{Answer:} False
\\ \textit{Options:} A) True; B) False
\\ \textit{Note:} Statistical significance under low power is unreliable because it increases the risk of type I errors, leading to false conclusions about the presence of an effect, especially when the sample size is small.
\item \textbf{Answer:} False
\\ \textit{Options:} A) True; B) False
\\ \textit{Note:} The statement is false because a psychometrician who accurately estimates IQ scores is demonstrating reliability and validity in their assessments, rather than a constant error, which would imply a systematic bias in their measurements.
\item \textbf{Answer:} True
\\ \textit{Options:} A) True; B) False
\\ \textit{Note:} Individuals with moderate mental retardation typically demonstrate learning capabilities that allow them to achieve academic skills equivalent to about the eighth-grade level, often through specialized educational interventions and support systems tailored to their needs.
\item \textbf{Answer:} True
\\ \textit{Options:} A) True; B) False
\\ \textit{Note:} The difference threshold, also known as the just noticeable difference (JND), is defined as the smallest amount of change in a stimulus that can be detected by an individual 50\% of the time, confirming that the statement is true.
\end{enumerate}

\section*{Long Form Response}
\begin{enumerate}[label=\arabic*.]
\setcounter{enumi}{10}
\item \textbf{Answer:} Sherman's and Key's (1932) study of the Hollow Children critically examined the impact of environmental factors on intellectual development, yet ultimately revealed no significant findings. This lack of evidence suggests that the anticipated correlations between environmental influences-such as socio-economic status, educational opportunities, and familial support-and cognitive outcomes may not be as straightforward as previously thought. The study highlights the complexity of intellectual development, indicating that while environment plays a role, it may not be the sole determinant of cognitive abilities. Consequently, these findings prompt further investigation into other contributing factors, such as genetic predispositions and individual differences, that may interact with environmental elements to shape intellectual growth.
\\ \textit{Note:} correct because it accurately summarizes Sherman's and Key's findings, emphasizing the unexpected lack of significant correlations between environmental factors and intellectual development, thereby illustrating the complexity of this relationship.
\item \textbf{Answer:} Generalization in classical conditioning refers to the phenomenon where a conditioned response is elicited not only by the conditioned stimulus but also by stimuli that are similar to it. In Ivan Pavlov's experiments with dogs, after the dogs learned to associate the sound of a bell (the conditioned stimulus) with food (the unconditioned stimulus), they began to salivate in response to other similar sounds, such as a buzzer or a tuning fork. This illustrates broader principles in learning theory by demonstrating how organisms can adapt their responses based on similarities in their environment, indicating that learning is not just about specific stimuli but also about the general principles that govern behavior. Generalization emphasizes the efficiency of learning processes, allowing organisms to respond appropriately to a range of stimuli without needing to form new associations for each one.
\\ \textit{Note:} correct because it accurately describes generalization in classical conditioning, exemplified by Pavlov's dogs salivating to stimuli similar to the original conditioned stimulus, thereby demonstrating a key principle of how learning can transfer across similar experiences.
\end{enumerate}

\vfill
\noindent\textit{End of Answer Key}
\end{document}